

\fsection[\textcolor{waveBlue2}{ejercicio 6}]{\boldit{\textcolor{waveBlue2}{Investiga}}{
Busca tres ejemplos de empresas cuya actividad se base en la
economía circular . Explica a qué se dedican , cuál es su modelo de negocio y
cuáles de las 7R de la economía circular aplican .
}}

\begin{solution}

\textbf{1. ECOALF (España)}
\begin{itemize}
    \item \textbf{Actividad:} Moda sostenible; ropa, calzado y accesorios.
    \item \textbf{Modelo de negocio:} Recogen residuos como botellas de plástico, redes de pesca y basura marina, los reciclan y convierten en tejidos de calidad. Venden los productos al consumidor final.
    \item \textbf{7R aplicadas:}
        \begin{itemize}
            \item Reciclar: convierten residuos en nuevas materias primas.
            \item Reducir: disminuyen consumo de recursos naturales.
            \item Reutilizar: los residuos tienen una segunda vida en los productos.
        \end{itemize}
\end{itemize}

\textbf{2. Patagonia (EE. UU.)}
\begin{itemize}
    \item \textbf{Actividad:} Ropa y equipo outdoor sostenible.
    \item \textbf{Modelo de negocio:} Promueven reparación y reventa de ropa usada mediante el programa “Worn Wear”; venden productos duraderos y reciclables.
    \item \textbf{7R aplicadas:}
        \begin{itemize}
            \item Reparar: fomentan reparación de productos.
            \item Reutilizar: reventa de ropa usada.
            \item Reciclar: materiales reciclados en nuevos productos.
        \end{itemize}
\end{itemize}

\textbf{3. Loop Industries (Canadá)}
\begin{itemize}
    \item \textbf{Actividad:} Reciclaje químico de plásticos PET.
    \item \textbf{Modelo de negocio:} Recogen plásticos de desecho, los transforman químicamente en PET de alta calidad para fabricar botellas y envases, que venden a fabricantes de envases.
    \item \textbf{7R aplicadas:}
        \begin{itemize}
            \item Reciclar: plásticos reciclados químicamente.
            \item Reducir: disminuyen necesidad de PET virgen.
            \item Reintroducir: los materiales reciclados vuelven al ciclo productivo.
        \end{itemize}
\end{itemize}
\end{solution}
